\subsection{绝对半单模}

定义 1. —— 设 K 是一个交换域，A 是一个 K-代数。若对 K 的每个扩张 L，\(\mathrm{A}_{(\mathrm{L})}\) -模 \(\mathrm{M}_{(\mathrm{L})}\) 都是半单的，则称 A-模 M 是绝对半单的。

每个绝对半单模都是半单的。反之，若域 K 是完满的，则每个在 K 上有限维的半单 A-模都是绝对半单的（VIII, 第 222 页，推论 1, a)），因为完满域的每个扩张都是可分的（V, §7, No.~1, 第 36 页，命题 2）。

命题 1. —— 设 K 是一个交换域，A 是一个 K-代数。

\begin{enumerate}
\def\labelenumi{\alph{enumi})}
\item
  绝对半单 A-模的任意直和是一个绝对半单 A-模。绝对半单模的每个子模与商模都是绝对半单的。
\item
  设 M 是一个 A-模，L 是域 K 的一个扩张。A-模 M 是绝对半单的当且仅当 \(\mathrm{A}_{(\mathrm{L})}\) -模 \(\mathrm{M}_{(\mathrm{L})}\) 是绝对半单的。
\end{enumerate}

断言 a) 由关于半单模的类似断言得出（VIII, 第 56 页，推论 1 与 3）。

假设 A-模 M 是绝对半单的，并设 \(L'\) 是 L 的一个扩张。作为 \(\mathrm{A}_{(\mathrm{L}')}\) -模，\(L' \otimes_{L} M_{(L)}\) 同构于 \(\mathrm{M}_{(\mathrm{L}')}\)（II, §5, No.~1, 第 273 页，命题 2）；因此它是一个半单模。这就证明了 \(\mathrm{M}_{(\mathrm{L})}\) 是绝对半单的。

反之，假设 \(M_{(L)}\) 是绝对半单的。设 \(L'\) 是 K 的一个扩张。存在一个复合扩张 \((\Omega, u, v)\)（L 与 \(L'\) 的；V, §2, No.~4, 第 13 页，推论）；我们把 L 与 \(L'\) 等同为 \(\Omega\) 的子扩张。

\(A_{(\Omega)}\) -模 \(M_{(\Omega)}\) 同构于 \((\mathrm{M}_{(\mathrm{L})})_{(\Omega)}\)；因此它是半单的。但 \(M_{(\Omega)}\) 也同构于 \((\mathrm{M}_{(\mathrm{L}^{\prime})})_{(\Omega)}\)，且 VIII, 第 222 页的命题 8, a) 蕴含 \(M_{(L^{\prime})}\) 是半单的。于是 M 是绝对半单的。

命题 2. —— 设 K 是一个交换域，A 是一个 K-代数，M 是一个在 K 上有限维的 A-模。下列性质等价：

\begin{enumerate}
\def\labelenumi{(\roman{enumi})}
\item
  A-模 M 是绝对半单的。
\item
  存在 \(\mathrm{P}\)（\(\mathrm{K}\) 的一个为完满域的扩张），使得 \(\mathrm{A}_{(\mathrm{P})}\) -模 \(\mathrm{M}_{(\mathrm{P})}\) 是半单的。
\item
  A-模 M 是半单的，且其交换子的中心是域 K 上的一个平展代数。
\end{enumerate}

显然 (i) 蕴含 (ii)。假设 (ii) 成立；则由 VIII, 第 222 页的推论 1, a)，\(A_{(P)}\) -模 \(M_{(P)}\) 是绝对半单的。由命题 1, b)，M 于是是绝对半单的。因此 (ii) 蕴含 (i)。

假设 M 是半单的。设 Z 是 M 的交换子的中心；它是域 K 上的一个有限次数交换代数。若 L 是 K 的一个扩张，则由 VIII, 第 222 页的命题 8, b) 得出，\(A_{(L)}\) -模 \(M_{(L)}\) 是半单的当且仅当环 \(Z_{(L)}\) 是约化的。因此 (i) 与 (iii) 的等价性由 V, §6, No.~7, 第 34 页的定理 4 得出。

